\section{BIRD Execution Traces and Supporting Analyses}\label{app:bird-traces}
This appendix expands the supporting execution diagnosis. In the
historical records, Phase I denotes BIRD discovery and Phase II the
separate BIRD admission study; neither explains the MATH results causally.
\subsection{Discovery populations and conditions}\label{app:phase1-provenance}

\paragraph{Population and judging.}
The discovery comparison in \S\ref{sec:f1-collapse} evaluates twelve
TTHE-proposer candidates from a GLM-5.3-Flash builder, plus \bare{}
and \react{}, on $60$ tasks. Admission checks import and smoke execution;
it does not validate the claimed mechanism. The baseline flipped three
of the $60$ items between discovery runs, so their baseline-error sets
differ slightly. Discovery used a set-of-rows proxy judge. On the two
matrices with retained final SQL, its agreement with the official BIRD
scorer is 0.97--0.99, with no proxy-correct/official-incorrect cases.

\paragraph{Development exposure and record limits.}
The hand-written controls were designed after inspecting discovery
traces; $24$ of their $151$ evaluation items overlapped
protocol-development measurements. Their headroom is descriptive,
not a blind estimate of stable complementarity. Early in-session
generation records do not retain the generation tool and per-attempt
logs needed to audit equal generation budgets. These limitations
apply to the discovery evidence retained here.

\begin{figure}[H]
\centering
\includegraphics[width=0.60\linewidth]{figures/fig_phenomenon.pdf}
\caption{\textbf{Source differences need not change observed correctness.}
BIRD discovery: $14$ harnesses, $60$ tasks, one execution each.
Across $91$ pairs including baselines, token-sequence source similarity
has $\rho=-0.010$ rank association with correctness disagreement;
$36/91$ pairs have identical correctness vectors. This is an outcome
comparison on the observed tasks, not a test of program equivalence.}
\label{fig:phenomenon}
\end{figure}
\FloatBarrier
\subsection{Discovery measurements and hand-written controls}\label{app:bird-discovery}

Every BIRD discovery candidate passes a code-difference check. Yet
$21/66$ generated-candidate pairs have identical correctness vectors
($36/91$ including baselines), and source similarity has essentially
no rank association with correctness disagreement across the $91$
pairs ($\rho=-0.010$; Figure~\ref{fig:phenomenon}). Mean disagreement
is $3.3\%$ overall ($3.8\%$ among generated candidates), union repair
of \bare{} errors is $7.7\%$, and one-run headroom is $3.3$\,pp.
Agreement on these $60$ tasks does not establish general program
equivalence, but it motivates examining what the programs actually do.


\subsection{Traces distinguish absence, non-activation, and failure}
\textbf{T1: described but absent.} Six of twelve candidates claim
retrieval or self-checking in docstrings but make a single model call
and execute no artifact. The named mechanism is not implemented.
\textbf{T2: present but untriggered.} In all ten audited \react{} tasks,
the first SQL executes successfully, so an existing
execute--feedback--repair branch never fires; final SQL equals
\bare{}. Successful execution is not semantic correctness, and these
runs do not test repair under failure.
\textbf{T3: executed but broken.} One of twelve candidates performs
multi-turn execution and repair but returns prose as final SQL,
losing tasks that \bare{} answers correctly.

These observations locate different failures rather than assigning
one diversity score. T1 requires checking implementation; T2 requires
tasks that exercise the branch; T3 requires a valid output contract.
They do not establish that activation causes improved accuracy.
Four hand-written controls with distinct paths produce $6.67$\,pp
one-run headroom on $60$ tasks and $7.28$\,pp on $151$ items, showing
that the apparatus can record different repairs. But they were written
after trace inspection and lack matched repeats and clone controls.
Appendix~\ref{app:phase1-provenance} discloses the legacy proxy and
$24/151$ items overlapping protocol development;
Appendix~\ref{app:examples} gives worked traces.


\FloatBarrier
\subsection{Worked examples from the trace audit}\label{app:examples}

The T1 count above classifies all twelve discovery candidates. The separate
trace audit covers five harnesses: \bare{}, \react{}, the single-call
candidate, the timeout-retry candidate, and the multi-turn repair candidate.
It illustrates execution paths rather than providing the denominator for
the six-of-twelve T1 count. This section walks one task through those
paths. Outputs are quoted from this audit (5 harnesses $\times$ 10 tasks); its logger caps
stored SQL excerpts at 120 characters, so truncations are marked
\emph{[trunc.]}.

\paragraph{The task.} BIRD dev-set question 352 (\texttt{card\_games}
database; question text verbatim, the typo is the dataset's):
\begin{quote}\small\ttfamily
Calculate the percentage of the cards availabe in Chinese Simplified.
\end{quote}
Its BIRD hint:
\begin{quote}\small\ttfamily
'Chinese Simplified' is the language; percentage = Divide(Sum(id where
language = 'Chinese Simplified'), Count(id)) *100
\end{quote}
Gold SQL (executing the two on the released database yields 8.7734 and
8.7728 respectively --- the official scorer separates them):
\begin{quote}\small\ttfamily
SELECT CAST(SUM(CASE WHEN T2.language = 'Chinese Simplified' THEN 1 ELSE 0
END) AS REAL) * 100 / COUNT(T1.id) FROM cards AS T1 INNER JOIN foreign\_data
AS T2 ON T1.uuid = T2.uuid
\end{quote}
The \bare{} harness answers:
\begin{quote}\small\ttfamily
SELECT SUM(CASE WHEN language = 'Chinese Simplified' THEN 1 ELSE 0 END) *
100.0 / COUNT(id) AS percentage FROM foreign\_data;
\end{quote}
marked wrong by the official scorer: it counts directly over
\texttt{foreign\_data}, whereas the gold query counts after joining to
\texttt{cards}; the hint does not spell out that join.

\paragraph{\react{} (T2: mechanism present but untriggered).} \react's control
flow does execute SQL and retries on database errors. But its first call is
byte-identical to \bare's prompt, and here the generated SQL executes without
error, so the error-triggered retry never fires and the final SQL equals
\bare's. Across all 10 audited tasks the pattern held: SQL executions
present, final SQL byte-identical to \bare's on 10/10 --- the feedback loop
is real code whose repair branch did not activate on these tasks. Controlled
failure injection is needed to test whether feedback changes the output.

\paragraph{Single-call candidate (T1: mechanism described, not
executed).} This candidate's docstring claims a ``conservative repair'':
cap the solver's maximum output tokens and demand an answer-only block.
Its behavior is structurally \bare's: one LLM call followed by
SQL extraction. The token cap exists only as an environment default (2{,}048 against the runtime's
32{,}000) --- no harness control flow implements it, and it never binds on
these tasks because a SQL answer is orders of magnitude below the cap. On
the worked task it issues exactly one call, executes nothing, and returns
the same wrong query as \bare{} modulo whitespace (counting over
\texttt{foreign\_data} without the join). Across the 10 audited tasks it
records zero SQL executions on all 10, and its final SQL is byte-identical
to \bare's on 6 (the remaining differences trace to the reworded system
prompt, with no implemented retrieval, self-check, or repair operation).
Prompt rewriting is an implemented transformation; it does not implement
the additional mechanism claimed in the docstring. The most common
Phase-I failure is thus a claimed operation absent from the executed program.

\paragraph{Timeout-retry candidate (code-level diversity,
behavioral identity).} This candidate is a genuinely different program: it
wraps the single coder call in a daemon-thread watchdog and fires one hedged
duplicate re-ask if the call stalls past 40\,s. The mechanism exists in the
control flow --- but its trigger (a stalled API call) never occurred in the
audit: zero SQL executions, and the final SQL equaled \bare's on all 10
audited tasks. Code difference without behavioral difference survives even
when the added mechanism is real: it simply never fires.

\paragraph{Multi-turn repair candidate (T3: a broken multi-turn
path).} The one audited candidate implementing genuine execute-inspect-repair
returns the model's \emph{prose} as final SQL. On a task \bare{} answered
correctly (``State the alternative languages available for card named Annul
numbered 29.''), its logged final SQL begins, verbatim:
\begin{quote}\small\ttfamily
with "only a single read-only SELECT is allowed". The previous "query" was
actually just prose text (the model wrote rea\emph{[trunc.]})
\end{quote}
The repair re-prompt invites the failure: asked to fix a query that ``was
actually just prose'', the frozen target responds with more prose reasoning,
no SQL block is found, and the reasoning text is submitted as the query. The
item flips from single-call-correct to harness-wrong.


\FloatBarrier
